% !PRINTED 536
\weberpage{536}
%<<CONT
which likewise has the discriminant $-4m$, and in which, if $m \equiv 1$
or $\equiv 2\ (\mathrm{mod}\ 4)$, the two outer coefficients are odd; for
we have:
\[
\begin{aligned}
&\text{for } m \equiv 1\ (\mathrm{mod}\ 4), \quad B \equiv 0\ (\mathrm{mod}\ 2), \quad A+C \equiv 2\ (\mathrm{mod}\ 4),\\
&\text{for } m \equiv 2\ (\mathrm{mod}\ 4), \quad B \equiv 1\ (\mathrm{mod}\ 2), \quad A+C \equiv 0\ (\mathrm{mod}\ 4).
\end{aligned}
\]

From this it follows that
\[
f(\omega) \quad\text{and}\quad f\left(\dfrac{\omega-1}{\omega+1}\right),
\]
apart from 24th roots of unity, run through the same sequence of values.
For if one replaces $\omega'$ by an equivalent number, then, if the outer
coefficients are to remain odd, the substitution must belong to the first
or the second class (§ 36), and from (4) it then follows that $\omega$ too
passes into an equivalent number. But if two forms (5) are equivalent,
then the corresponding forms (3) are equivalent as well, and conversely.

Accordingly
\[
\Pi\, f(\omega)\, f\!\left(\frac{\omega-1}{\omega+1}\right) = 2^{2h\tau},
\]
while on the other hand [§ 34, (18)]:
\[
f(\omega)\, f\!\left(\frac{\omega-1}{\omega+1}\right) = \sqrt{2},
\]
from which there results:
\weq{6}{\tau = \frac{1}{4}, \quad m \equiv 1,\, 2\ (\mathrm{mod}\ 4),}
in agreement with the result of the preceding section [§ 141, (31)].

2. If next $m \equiv 3\ (\mathrm{mod}\ 8)$, then to a root
$\omega'$ of a form of discriminant $-m$ there correspond in each case three
roots of forms of discriminant $-4m$:
\[
2\omega', \quad \frac{\omega'}{2}, \quad \frac{\omega'+1}{2},
\]
and the 24th powers of the quantities:
\[
\begin{aligned}
f_1(2\omega') &= \frac{\sqrt{2}}{f_2(\omega')},\\
f_2\!\left(\frac{\omega'}{2}\right) &= \frac{\sqrt{2}}{f_1(\omega')},\\
f_2\!\left(\frac{\omega'+1}{2}\right) &= \frac{\sqrt{2}}{f(\omega')},
\end{aligned}
\]
%CONT>>
